\section{Methods}

\subsection{Connectomes, and what each one can support}
\label{sec:data}

We analyse four public connectome releases spanning three phyla, from which we take
five networks --- the \textit{Drosophila} release supplies two functionally distinct
circuits (Table~\ref{tab:datasets}). Two of the releases carry synaptic valence and two
do not, and this distinction constrains which questions each dataset can answer; we
state it here rather than in a caption because it is the reason several rows of our
tables are deliberately empty.

\textit{Caenorhabditis elegans.} The adult hermaphrodite somatic connectome
\cite{cook2019whole,varshney2011structural}, with neurotransmitter assignments
giving an excitatory/inhibitory (E/I) label per neuron.

\textit{Drosophila melanogaster.} FlyWire v630 \cite{dorkenwald2024neuronal,
lin2024network}, from which we take the central complex and the antennal lobe as
two functionally distinct circuits. Valence follows the released neurotransmitter
predictions.

\textit{Mus musculus.} The MICrONS Pinky100 V1 layer 2/3 proofread subgraph
\cite{microns2025functional}. This release annotates 34 inhibitory somata, but every
one of them is disconnected in the accompanying synapse table: all 1{,}736 synapses
are excitatory$\rightarrow$excitatory. We therefore treat this network as
\textbf{topology-only} and do not compute sign-attributed motifs on it. Reporting
coherent feedforward loops here would be reporting plain feedforward loops under
another name, since with no inhibitory neuron in the graph the two counts are
identically equal.

\textit{Homo sapiens.} The Human Connectome Project group-consensus structural
network \cite{vanessen2012human}, 1{,}015 parcellated regions from diffusion
tractography. Streamline density carries no synaptic valence, so this network is
likewise treated as topology-only. We emphasise that it is a regional network, not a
neuronal connectome, and we make no microscale claim from it.

We do not analyse human microscale connectivity. The H01 release
\cite{shapsoncoe2024petavoxel} provides cell and synapse tables but the connectivity
export needed to build a wiring diagram is not part of the distribution we worked
from, and a census of a synapse table is not a connectome analysis.

Reservoir capacity grows with neuron count, so cross-species comparison requires
size matching. From each connectome we extract an induced subgraph of at most
$N = \NnodesTarget$ neurons by breadth-first expansion from a seed neuron of
above-median degree, following both afferent and efferent edges; degree-zero nodes
are then removed. Uniform node sampling is not used, because it destroys degree
structure and motif density by construction.

\subsection{Exact sign-attributed motif census}
\label{sec:census}

For a directed graph with neuron valences $\sigma_i \in \{+1,-1\}$ and Dale's
principle -- every synapse carries the valence of its presynaptic neuron -- the sign
product around a feedforward loop $(u \to v,\; v \to w,\; u \to w)$ is
\begin{equation}
\sigma_u \cdot \sigma_v \cdot \sigma_u \;=\; \sigma_u^2\,\sigma_v \;=\; \sigma_v .
\label{eq:coherence}
\end{equation}

\begin{remark}
A feedforward loop is coherent if and only if its \emph{intermediate} neuron is
excitatory. Coherence is a property of one node, not of a product over three edges.
\end{remark}

Equation~\eqref{eq:coherence} makes the sign-attributed census exactly as cheap as
the topological one and removes any need to store per-edge signs. We count
\begin{align}
\mathrm{FFL} &= \sum_{(u,v) \in E} \big| N^{+}(u) \cap N^{+}(v) \big| ,\\
\mathrm{cFFL} &= \sum_{\substack{(u,v)\in E \\ \sigma_v = +1}} \big| N^{+}(u) \cap N^{+}(v) \big| ,
\end{align}
with $\mathrm{iFFL} = \mathrm{FFL} - \mathrm{cFFL}$; reciprocal pairs, mutually
inhibitory pairs and directed 3-cycles follow analogously, and bi-fans are counted
exactly by accumulating source-pair co-occurrence through shared targets rather than
by sampling node pairs. All counts are exact; none are estimated. Single-edge
insertions and deletions admit exact $O(d)$ updates of every 3-node count, which is
what makes the annealed null model of \S\ref{sec:nulls} tractable.

\subsection{Null ensembles and the structural control ladder}
\label{sec:nulls}

\paragraph{Dale's principle constrains nothing here.}
The conventional null is the Maslov--Sneppen double-edge swap
\cite{maslov2002specificity}, which fixes the joint in/out degree sequence. It is
common to describe a connectome null as preserving degrees \emph{and} Dale's
principle. Under the convention above this is not an additional constraint: the swap
$(u\to v,\, x\to y) \Rightarrow (u\to y,\, x\to v)$ leaves $u$ and $x$ presynaptic,
so every neuron's outgoing synapses retain a single shared valence automatically.
A ``Dale-compliant'' rewiring is therefore the standard ensemble, and it cannot test
whether the \emph{placement} of inhibition in the graph is non-random. That requires
permuting valence labels across neurons, which we include as a separate rung.

\paragraph{Two weighted rewirings, and which one we report.}
For the weighted \texttt{conn2res} connectomes the swap above underdetermines the
ensemble, because it says what happens to edges and not what happens to weights. We
implement both readings and report both (\S\ref{sec:res-nullchoice}).

\emph{Weight-carrying rewiring} moves each weight with the edge that carries it, and
rejects swaps that disconnect the graph. We obtain it by calling
\texttt{bct.reference.randmio\_und\_connected} or \texttt{randmio\_dir\_connected}
\cite{rubinov2010complex} directly, dispatching on symmetry and using the toolbox default
of ten swap attempts per edge, rather than reimplementing it --- the point of
\S\ref{sec:res-nullchoice} is that a reimplementation from a description is not
guaranteed to be the same ensemble. This is what \texttt{conn2res} calls, so it is the
ensemble every comparison in this paper reports unless stated otherwise.

\emph{Rank-reassigned rewiring} randomises the binary topology by the same swaps and then
assigns the original multiset of weight values to the new edges in decreasing order of
endpoint-degree product $d^{\mathrm{out}}_u d^{\mathrm{in}}_v$. It preserves the degree
sequence and the weight distribution exactly and destroys the association between a
weight and its edge. Both constructions satisfy the description ``randomised while
preserving density and degree sequence''; they are different distributions over graphs;
and \S\ref{sec:res-nullchoice} reports what that difference is worth.

Both are applied \emph{after} the toolbox's preprocessing (min--max scaling followed by
division by the spectral radius) and before the gain is set, so that each arm of a
comparison is a randomisation of the same object.

\paragraph{The motif-preserving ensemble.}
The rung the connectome literature does not use is one that preserves the motif census
itself. Constructions for it exist outside that literature and we do not claim the idea:
Milo et al.\ proposed a simulated-annealing null preserving lower-order subgraph counts
\cite{milo2002network,milo2004superfamilies}; exponential random graph models condition
on triad counts by design \cite{frank1986markov,robins2007introduction}; the $dk$-series
is a nested hierarchy of nulls constraining progressively higher-order structure, with
$2K$ and $3K$ fixing degree correlations and connected triads
\cite{mahadevan2006systematic,orsini2015quantifying}; and Fischer et al.\ sample
motif-constrained ensembles at fixed degree sequence by Wang--Landau multicanonical
sampling, which unlike annealing comes with a sampling guarantee
\cite{fischer2015sampling}. What we contribute is an implementation for the
\emph{sign-attributed} census at connectome scale --- exact coherent and incoherent
feedforward-loop counts under Dale's principle, with $O(d)$ incremental updates and a
move set that survives near-total reciprocity --- and the use of that ensemble as a
control rung rather than as an enrichment null. Where a formally exact sampler is
available it should be preferred, and \S\ref{sec:res-robust} reports what our annealed
chain can and cannot demonstrate in its place. We sample it with a Metropolis chain over degree-preserving swaps whose
energy is the squared deviation of the running census from a target census $c^{*}$,
measured in units of the unconstrained-ensemble standard deviation $s_m$:
\begin{equation}
\mathcal{E}(G) \;=\; \sum_{m \in \mathcal{M}} \left( \frac{c_m(G) - c^{*}_m}{s_m} \right)^{2},
\label{eq:anneal}
\end{equation}
over $\mathcal{M} = \{$reciprocal pairs, mutually inhibitory pairs, FFL, cFFL, iFFL,
3-cycles$\}$, over $\SweepCount \times |E|$ proposals with a geometric temperature
schedule. Exact incremental updates (\S\ref{sec:census}) make each proposal $O(d)$.
Setting $c^{*}$ to the empirical census yields graphs that are maximally randomised
subject to reproducing the connectome's motif spectrum.

Two things about the construction are worth stating precisely, because the obvious
implementation of it does not work. First, where the chain starts determines which
property is guaranteed and which must be demonstrated. Starting from a fully rewired
graph and annealing towards $c^{*}$ guarantees randomisation and leaves convergence to
be shown; starting from the connectome itself and running a low-temperature constrained
walk guarantees convergence --- the chain begins on the constraint surface --- and
leaves randomisation to be shown. These are equivalent where both succeed, but the
random start fails outright on networks whose census is hard to rebuild by chance. We
therefore use the empirical start and demonstrate randomisation directly, by edge
overlap against the empirical graph and against an unconstrained rewiring
(\S\ref{sec:res-null}).

Second, an ordinary double-edge swap destroys reciprocal pairs far more readily than it
creates them. In a near-symmetric network --- diffusion tractography above all, where
almost every edge is reciprocal --- the reciprocal-pair count therefore drifts down by
tens of percent however stiff the penalty, and the constrained ensemble is effectively
unreachable by that move set. We add a move that swaps two reciprocal \emph{pairs} as
units, $(u \leftrightarrow v,\, x \leftrightarrow y) \mapsto (u \leftrightarrow y,\,
x \leftrightarrow v)$, which preserves every node's in- and out-degree and the
reciprocity of all four edges, and mix it in proportion to the fraction of edges that
are reciprocal.
Bi-fan is deliberately excluded from $\mathcal{M}$ and reported post hoc as a
held-out statistic, which tests whether the chain reproduces structure it was not
asked to reproduce.

\paragraph{The ladder.}
Each rung adds exactly one structural property to the one below it, so the
contribution of that property is the difference between adjacent rungs:
\begin{description}[leftmargin=1.1em,itemsep=1pt]
  \item[L0] Erd\H{o}s--R\'enyi: node count, edge count and E:I ratio only.
  \item[L1] $+$ joint in/out degree sequence (valence labels permuted across nodes).
  \item[L2] $+$ valence placement (labels left on their own neurons).
  \item[L3] $+$ 3-node motif census (the annealed ensemble above).
  \item[L4] the empirical connectome.
\end{description}
Every rung has identical node and edge counts by construction. We draw
$\NinstancesLadder$ independent surrogates per rung. Rung L4 is different in kind: the
connectome is one graph, so its $\NinstancesLadder$ replicates differ in the input
coding and the sensory/motor split only. Its replicates therefore carry one variance
component where L0--L3 carry two, and we state the consequence rather than leave it to
be inferred. Contrasts use Welch's test, which does not assume equal variances, but a
contrast involving L4 supports inference about \emph{this} connectome rather than about
connectomes as a class; the surrogate rungs support the stronger reading. The asymmetry
works against us where our conclusion is a null one: in the L3--L4 contrast a real
difference specific to the empirical graph would be easier, not harder, to detect at
reduced variance, so a small measured effect there is not an artefact of the variance
structure.

\subsection{Task and evaluation protocol}
\label{sec:task}

\paragraph{Benchmark.}
We use the Spiking Heidelberg Digits (SHD) dataset \cite{cramer2022heidelberg}:
20 spoken-digit classes in English and German, converted to spike trains by a
cochlear model, from 12 speakers, with a canonical test split that holds out two
speakers entirely. Spikes are binned into $\Tsteps$ time steps over 1\,s and the 700
cochlear channels are pooled into $\Cchannels$ bands. We use a real spike-based
benchmark rather than a synthetic pattern task because synthetic tasks of the kind
used in earlier work saturate near 100\% accuracy for every architecture and
therefore cannot discriminate between them.

\paragraph{Reservoir protocol.}
The connectome or surrogate supplies the recurrent weight matrix and is never
trained; only a ridge readout is fitted \cite{lukosevicius2009reservoir,
suarez2024connectome}. Neurons are discrete-time leaky integrate-and-fire units,
\begin{equation}
v_i(t{+}1) = \lambda\, v_i(t)\,\big(1 - s_i(t)\big) + I_i(t),
\qquad
s_i(t) = \mathbb{I}\!\left[v_i(t) \ge \vartheta\right],
\end{equation}
with $I(t) = x(t) W_{\mathrm{in}} + s(t) W_{\mathrm{rec}}$, leak $\lambda =
\Leak$ and threshold $\vartheta = 1$.

\paragraph{Controls that the comparison depends on.}
Four quantities are held fixed across every network in a comparison: node count, edge
count, spectral radius (each $W_{\mathrm{rec}}$ is rescaled to $\rho =
\SpectralRadius$), and the input projection. Without spectral-radius matching the
comparison measures effective gain rather than topology.

A fifth control is less obvious and turned out to be decisive. If every neuron
receives the stimulus and every neuron feeds the readout, then the task is solved by
the input projection and the per-neuron leak alone: under that protocol a reservoir
with $W_{\mathrm{rec}} = 0$ scores $\ZeroRecFull$ on SHD, statistically
indistinguishable from any connectome, and no structural comparison can detect
anything. We therefore drive a sensory subset ($\FracInput$ of neurons) and decode
from a \emph{disjoint} motor subset ($\FracReadout$), so task-relevant signal must
traverse the recurrent graph to be decodable. The zero-matrix control then falls to
$\ZeroRecSplit$ against a chance level of $5.0\%$, which is the diagnostic that the
benchmark is sensitive to recurrent structure at all. We recommend that any
structural comparison of this kind report the corresponding control.

\paragraph{Matching the operating point, not only the gain.}
A sixth control is needed because spectral radius fixes the recurrent gain but not the
activity level the network settles into. At the common $\rho = \SpectralRadius$ the
rungs do not fire at the same rate: the Erd\H{o}s--R\'enyi surrogates run at roughly
twice the mean rate of the connectomes they are matched to (\S\ref{sec:res-rate}), and
firing rate is both correlated with accuracy across our runs ($r = \RateAccCorr$) and
the direct determinant of the synaptic-operation count used as an energy proxy. A
contrast at matched $\rho$ therefore confounds topology with operating point, and in
the direction that flatters the unstructured control.

We consequently report the ladder twice. In the \emph{matched-gain} version every
network is set to $\rho = \SpectralRadius$, and the rung differences include whatever
change in activity level the topology induces. In the \emph{matched-rate} version each
network is calibrated individually, by bisection on its own spectral radius, to a common
mean firing rate of $\TargetRate$, so that the comparison is made at a common operating
point and the residual difference is attributable to wiring rather than to activity. The
two versions answer different questions --- at equal gain, what operating point does
each topology select; and at equal operating point, does topology still change what the
network computes --- and we regard reporting only the first as the more common error.

\paragraph{Readout and statistics.}
Features are the time-averaged firing rate and the final membrane potential of the
motor population. A ridge readout is fitted with the regulariser selected on a
held-out validation split, never on test. Energy is reported as synaptic operations,
$\mathcal{S} = \sum_t \sum_{i: s_i(t)=1} k^{\mathrm{out}}_i$. Comparisons between
rungs use Welch's $t$-test with Holm correction \cite{holm1979simple} across the full
family of contrasts, and are reported with percentile-bootstrap confidence intervals
and Hedges' $g$, because with $\NinstancesLadder$ instances per rung a detectable
difference can still be operationally negligible.

\subsection{Perturbation theory without a spectral gap}
\label{sec:theory}

Degeneracy is naturally posed as matrix perturbation, and the standard tool is the
Davis--Kahan $\sin\Theta$ theorem \cite{daviskahan1970rotation,stewart1990matrix},
which bounds the rotation of an invariant subspace by $\|E\|/\delta$ for a spectral
gap $\delta$. On real connectomes this is vacuous: the smallest measured gap is
$\MinSpectralGap$ (\S\ref{sec:res-degeneracy}), so the bound admits perturbations of
norm at most $\MaxDKAdmissible$ --- far smaller than the removal of a single synapse
--- while the networks in fact tolerate removal of a large fraction of their edges.
That ratio is a statement about eigenvector instability under near-degenerate
spectra, not about brains.

What must be stable is not an eigenvector but the finite-horizon input--output map.
For horizon $T$ that map is the sequence of Markov parameters $H_t = W_{\mathrm{out}}
A^{t} W_{\mathrm{in}}$, $t = 0,\dots,T-1$.

\begin{proposition}[Gap-free sensitivity of the transfer map]
\label{prop:gapfree}
Let $a = \|A\|_2$. For a perturbation $A \mapsto A + E$, to first order in $E$,
\begin{equation}
\|\Delta H_t\|_F \;\le\; t\,a^{\,t-1}\,\|E\|_F\,\|W_{\mathrm{in}}\|_2\,\|W_{\mathrm{out}}\|_2 ,
\end{equation}
and hence, aggregating over the horizon,
\begin{equation}
\Big(\textstyle\sum_{t<T} \|\Delta H_t\|_F^{2}\Big)^{1/2} \;\le\; C(T,a)\,\|E\|_F,
\end{equation}
\begin{equation}
C(T,a) = \|W_{\mathrm{in}}\|_2 \|W_{\mathrm{out}}\|_2
         \Big(\textstyle\sum_{t<T} t^{2}a^{2(t-1)}\Big)^{1/2}.
\end{equation}
No spectral gap is assumed, and $C$ is finite for every $A$.
\end{proposition}

\noindent\emph{Proof.} Expand $(A+E)^{t} - A^{t} = \sum_{s<t} A^{s} E A^{t-1-s} +
O(\|E\|^{2})$ and bound each of the $t$ terms submultiplicatively by
$\|A\|_2^{s}\,\|E\|_F\,\|A\|_2^{t-1-s} = a^{\,t-1}\|E\|_F$; multiply by the input and
output operators and collect. \hfill$\square$

\begin{remark}
The constant is stated in the operator norm and not in the spectral radius. Substituting
$\rho(A)$ for $\|A\|_2$ would require $\|A\|_2 \le \rho(A)$, which holds only for normal
matrices; connectome matrices are markedly non-normal, and we measure
$\|A\|_2/\rho(A)$ up to $2.1$ across our networks at matched $\rho$. With $\rho$ in
place of $a$ the constant is too small by up to that ratio raised to the horizon, and
the measured alignment $\gamma$ --- which is defined as the fraction of the worst case
a direction realises --- can exceed $1$. This matters because transient amplification
in a non-normal recurrent network is precisely the regime the bound is meant to cover.
\end{remark}

Proposition~\ref{prop:gapfree} is worst-case over the \emph{direction} of $E$. Real
remodelling is not adversarial, and the quantity that expresses this is the fraction
of the worst case a given direction realises.

\begin{definition}[Operational alignment]
\label{def:gamma}
\begin{equation}
\gamma(E) \;=\; \frac{\big(\sum_{t<T}\|\Delta H_t\|_F^{2}\big)^{1/2}}{C(T,a)\,\|E\|_F} \;\in\; [0,1].
\end{equation}
\end{definition}

\begin{proposition}[Tolerance scales inversely with alignment]
\label{prop:tolerance}
If task performance degrades once the aggregated transfer-map perturbation exceeds a
task-determined tolerance $\tau$, then the largest admissible perturbation magnitude
satisfies $\|E\|_F^{\max} \ge \tau / \big(\gamma(E)\,C(T,a)\big)$, i.e. it scales
as $1/\gamma$.
\end{proposition}

This is falsifiable in a way the usual degeneracy demonstration is not. Showing that
a connectome survives \emph{random} pruning establishes only that random directions
are benign. We therefore construct the maximally damaging perturbation of the same
cardinality by ranking synapses on their exact first-order influence on the
horizon-aggregated map,
\begin{multline}
\frac{\partial}{\partial A_{ij}} \sum_{t<T} \big\|W_{\mathrm{out}} A^{t} W_{\mathrm{in}}\big\|_F^{2} \\
= 2 \sum_{t<T} \sum_{s<t} \Big[ \big(W_{\mathrm{out}} A^{\,t-1-s}\big)^{\!\top} H_t
\big(A^{s} W_{\mathrm{in}}\big)^{\!\top} \Big]_{ij},
\end{multline}
evaluated by accumulating forward states $A^{s}W_{\mathrm{in}}$ and adjoint states
$W_{\mathrm{out}}A^{j}$ iteratively, so no matrix power is ever formed. If a
connectome tolerated adversarial and random pruning equally,
Proposition~\ref{prop:tolerance} would be false. Every pruned network is renormalised
to the original spectral radius before evaluation, so measured degradation reflects
lost structure rather than lost gain, and $\gamma$ is computed on exactly the
perturbation applied.

\subsection{Synthesis flow}
\label{sec:synth}

The accompanying Verilog is synthesised with Yosys \cite{wolf2013yosys} and mapped
with ABC \cite{ajaykumar2019abc} onto the Nangate45 Open Cell Library
\cite{nangate45} at the typical corner. We report gate count, flip-flop count and
mapped pre-layout cell area, and nothing else. We do not report die area, power,
energy per operation or a clock frequency: this flow performs no placement, routing,
clock-tree synthesis or parasitic extraction, switching power cannot be estimated
without activity factors and extracted parasitics, and the toolchain ships no static
timing analyser, so a critical path would have to be guessed rather than measured. Earlier revisions of
this project reported such figures for a commercial 28\,nm flow; no such flow was
run, and those figures are withdrawn (see the repository's audit record).
